\chapter{Statements\label{chap:Statements}}
Statements update storage elements and determine the flow of control in a subprogram.

\ExampleDef{Statements}
\listingref{Statements} shows some examples of statements.
\ASLListing{Examples of statements}{Statements}{\definitiontests/Statements.asl}

\ChapterOutline
\begin{itemize}
  \item \FormalRelationsRef{Statements} defines the formal relations used for statements;
  \item \secref{PassStatements} defines pass statements;
  \item \secref{AssignmentStatements} defines assignment statements;
  \item \secref{SetterAssignmentStatements} defines setter assignment statements;
  \item \secref{DeclarationStatements} defines declaration statements;
  \item \secref{DeclarationStatementsElidedParameter} defines declaration statements with an elided parameter;
  \item \secref{SequencingStatement} defines sequencing statements;
  \item \secref{CallStatements} defines call statements;
  \item \secref{ConditionalStatements} defines conditional statements;
  \item \secref{CaseStatements} defines case statements;
  \item \secref{AssertionStatements} defines assertion statements;
  \item \secref{WhileStatements} defines while statements;
  \item \secref{RepeatStatements} defines repeat statements;
  \item \secref{ForStatements} defines for-looping statements;
  \item \secref{ThrowStatements} defines throw statements;
  \item \secref{TryStatements} defines try statements;
  \item \secref{ReturnStatements} defines return statements;
  \item \secref{PrintStatements} defines print statements;
  \item \secref{UnreachableStatement} defines unreachable statements;
  \item \secref{PragmaStatements} defines pragma statements.
\end{itemize}

%%%%%%%%%%%%%%%%%%%%%%%%%%%%%%%%%%%%%%%%%%%%%%%%%%%%%%%%%%%%%%%%%%%%%%%%%%%%%%%%
\FormalRelationsDef{Statements}
\paragraph{Syntax:} Statements are grammatically derived from $\Nstmt$.

\paragraph{Abstract Syntax:} Statements are derived in the abstract syntax from $\stmt$
and generated by $\buildstmt$.
\RenderRelation{build_stmt}


\paragraph{Typing:} Statements are annotated by $\annotatestmt$.
\RenderRelation{annotate_stmt}

\paragraph{Semantics:} Statements are evaluated by $\evalstmt$.
\RenderRelation{eval_stmt}

%%%%%%%%%%%%%%%%%%%%%%%%%%%%%%%%%%%%%%%%%%%%%%%%%%%%%%%%%%%%%%%%%%%%%%%%%%%%%%%%
\section{Pass Statements\label{sec:PassStatements}}
\hypertarget{def-passstatementterm}{}

\ASLListing{A \texttt{pass} statement}{semantics-spass}{\semanticstests/SemanticsRule.SPass.asl}

\subsection{Syntax}
\begin{flalign*}
\Nstmt \derives \ & \Tpass \parsesep \Tsemicolon &
\end{flalign*}

\subsection{Abstract Syntax}
\RenderTypes[remove_hypertargets]{stmt_pass}


\ASTRuleDef{SPass}

\RenderProseAndFormally{build_stmt_SPass}

\subsection{Typing}
\TypingRuleDef{SPass}
\ExampleDef{Typing a Pass Statement}
Annotating the \passstatementterm{} in \listingref{semantics-spass} does
not modify the \staticenvironmentterm{}.


\RenderProseAndFormally{annotate_stmt_SPass}
\CodeSubsection{\SPassBegin}{\SPassEnd}{../Typing.ml}

\subsection{Semantics}
\SemanticsRuleDef{SPass}
\ExampleDef{Evaluation of Pass Statements}
In \listingref{semantics-spass}, \texttt{pass;} does not modify the environment.


\RenderProseAndFormally{eval_stmt_SPass}
\CodeSubsection{\EvalSPassBegin}{\EvalSPassEnd}{../Interpreter.ml}

\hypertarget{def-assignmentstatementterm}{}
\section{Assignment Statements\label{sec:AssignmentStatements}}
\subsection{Syntax}
\begin{flalign*}
\Nstmt \derives \ & \Nlexpr \parsesep \Teq \parsesep \Nexpr \parsesep \Tsemicolon &
\end{flalign*}

\subsection{Abstract Syntax}
\RenderTypes[remove_hypertargets]{stmt_assign}


\ASTRuleDef{SAssign}

\RenderProseAndFormally{build_stmt_SAssign}

\subsection{Typing}
\TypingRuleDef{SAssign}
\ExampleDef{Typing an Assignment Statement}
In \listingref{semantics-sassign}, the assignment \verb|x = 3;|
is well-typed.


\RenderProseAndFormally{annotate_stmt_SAssign}
\CodeSubsection{\SAssignBegin}{\SAssignEnd}{../Typing.ml}

\subsection{Semantics}
There are two rules for evaluating assignments:
\begin{itemize}
\item \SemanticsRuleRef{SAssignCall} handles assignments where the \rhsexpression{}
      is a \callexpressionterm{} and the left-hand-side expression is a tuple.
\item \SemanticsRuleRef{SAssign} handles all other assignments.
\end{itemize}

Although the sequential semantics of both types of statements is the same,
their concurrent semantics are different.
\SemanticsRuleRef{SAssignCall} generates a different execution graph ---
one where each value of the left-hand-side tuple depends on the \executiongraphterm{}
for the corresponding (that is, at the same position) value returned from the \callexpressionterm{}.
In contrast, the \executiongraphterm{} generated by \SemanticsRuleRef{SAssign}
creates dependencies between the entire \executiongraphterm{} for the evaluated right-hand-side
expression and the entire \executiongraphterm{} for the left-hand-side \executiongraphterm{}.

The rules for assignments first produce a value for the right-hand side expression
and then complete the update to the environment via an appropriate rule for evaluating the
\assignableexpression{} on the left-hand-side of the assignment.
The rules for updating the \assignableexpression{} handle variables, tuples, bitvectors, etc.

\SemanticsRuleDef{SAssignCall}
\ExampleDef{Evaluation of Multi-variable Assignment from Subprogram Calls}
In \listingref{assigncallsemantics}, given that the function call \texttt{f(1)} returns a triple of values ---
$\nvint(1)$, $\nvint(2)$, and $\nvint(3)$
(each with its own associated execution graph),
the statement \texttt{(a,b,-) = f(1)} assigns the value $\nvint(1)$ to the mutable variable \texttt{a},
$\nvint(2)$ to the mutable variable~\texttt{b}, and discards $\nvint(3)$.

\ASLListing{Assignment from a call expression.}{assigncallsemantics}{\semanticstests/SemanticsRule.SAssignCall.asl}


\RenderProseAndFormally{eval_stmt_SAssignCall}
\CodeSubsection{\EvalSAssignCallBegin}{\EvalSAssignCallEnd}{../Interpreter.ml}

\SemanticsRuleDef{SAssign}
\ExampleDef{Evaluation of Assignment Statements}
In \listingref{semantics-sassign},
\texttt{x = 3;} binds \texttt{x} to $\nvint(3)$ in the environment where \texttt{x} is bound to
$\nvint(42)$, and $\newenv$ is such that \texttt{x} is bound to $\nvint(3)$.
\ASLListing{Evaluating an assignment}{semantics-sassign}{\semanticstests/SemanticsRule.SAssign.asl}


\RenderProseAndFormally{eval_stmt_SAssign}
\CodeSubsection{\EvalSAssignBegin}{\EvalSAssignEnd}{../Interpreter.ml}

\SemanticsRuleDef{LexprIsVar}
\RenderRelation{lexpr_is_var}

In \ExampleRef{Evaluation of Multi-variable Assignment from Subprogram Calls},
the \assignableexpressions{} \verb|a|, \verb|b| are variables,
and \verb|-| is a discarded \assignableexpression{}.

\RenderProseAndFormally{lexpr_is_var}

\section{Setter Assignment Statements\label{sec:SetterAssignmentStatements}}
\subsection{Syntax}
\begin{flalign*}
\Nstmt \derives \
   & \Ncall \parsesep \Nsetteraccess \parsesep \Teq \parsesep \Nexpr \parsesep \Tsemicolon &\\
|\ & \Ncall \parsesep \Nsetteraccess \parsesep \Nslices \parsesep \Teq \parsesep \Nexpr \parsesep \Tsemicolon &\\
|\ & \Ncall \parsesep \Tdot \parsesep \Tlbracket \parsesep \Clisttwo{{\Tidentifier}} \parsesep \Trbracket \parsesep \Teq \parsesep \Nexpr \parsesep \Tsemicolon &\\
\Nsetteraccess \derives \
   & \emptysentence &\\
|\ & \Tdot \parsesep \Tidentifier \parsesep \Nsetteraccess &
\end{flalign*}

\SyntacticSugarDef{SetterFieldAssignment}

If a setter is followed by a slice or field access, then the assignment to that bitslice is
implemented using the following Read-Modify-Write (RMW) behaviour:
\begin{itemize}
\item invoking the getter of the same name as the setter, with the same actual
      arguments as the setter invocation, excluding the right-hand side of the assignment;
\item performing the assignment to the bitslice or field of the result of the
      getter invocation; and
\item invoking the setter to assign the resulting value.
\end{itemize}
This is implemented by desugaring via \ASTRuleRef{DesugarSetter}.
%
See \ExampleRef{A Setter Call as a Read-Modify-Write}.

\ASTRuleDef{MakeSetter}
\RenderRelation{make_setter}


\RenderProseAndFormally{make_setter}

\ASTRuleDef{BuildSetterAccess}
\RenderRelation{build_setter_access}


\RenderProseAndFormally{build_setter_access}

\ASTRuleDef{DesugarSetter}
\RenderRelation{desugar_setter}


\RenderProseAndFormally{desugar_setter}

\ASTRuleDef{DesugarSetterSetfields}
\RenderRelation{desugar_setter_set_fields}


\RenderProseAndFormally{desugar_setter_set_fields}

\ASTRuleDef{ReadModifyWrite}
\RenderRelation{read_modify_write}

\ExampleDef{A Setter Call as a Read-Modify-Write}
\listingref{ReadModifyWriteSlice} shows how a setter call followed by a slice is
\desugared.
\ASLListing{A Setter followed by a Slice}{ReadModifyWriteSlice}{\syntaxtests/ASTRule.ReadModifyWriteSlice.asl}

\listingref{ReadModifyWriteField} shows how a setter call followed by a field access is
\desugared.
\ASLListing{A Setter followed by a Field Access}{ReadModifyWriteField}{\syntaxtests/ASTRule.ReadModifyWriteField.asl}


\RenderProseAndFormally{read_modify_write}

\ASTRuleDef{SetterAssign}

\RenderProseAndFormally{build_stmt_SetterAssign}

\subsection{Typing and semantics}
As given by applying the relevant rules to the desugared AST.

\hypertarget{def-declarationstatementterm}{}
\section{Declaration Statements\label{sec:DeclarationStatements}}
\subsection{Syntax}
\begin{flalign*}
\Nstmt \derives \ & \Nlocaldeclkeyword \parsesep \Ndeclitem \parsesep \option{\Nasty} \parsesep \Teq \parsesep \Nexpr \parsesep \Tsemicolon &\\
|\ & \Tvar \parsesep \Ndeclitem \parsesep \Nasty \parsesep \Tsemicolon &\\
|\ & \Tvar \parsesep \Clisttwo{\Tidentifier} \parsesep \Nasty \parsesep \Tsemicolon &\\
\end{flalign*}

\subsection{Abstract Syntax}
\RenderTypes[remove_hypertargets]{stmt_decl}

\ASTRuleDef{SDecl}

\RenderProseAndFormally{build_stmt_SDecl}

\subsection{Typing}
\TypingRuleDef{SDecl}
\ExampleDef{Typing Declaration Statements}
\listingref{typing-sdecl} shows well-typed declaration statements.
\ASLListing{Typing declaration statements}{typing-sdecl}{\typingtests/TypingRule.SDecl.asl}

The specifications in \listingref{typing-sdecl-bad1} ill-typed,
since local constant storage elements are not allowed.
\ASLListing{A (illegal) local constant declaration}{typing-sdecl-bad1}{\typingtests/TypingRule.SDecl.bad1.asl}

The specification in \listingref{typing-sdecl-bad2} is ill-typed,
since immutable local storage elements require an initializing expression.
\ASLListing{An uninitialized immutable local storage declaration}{typing-sdecl-bad2}{\typingtests/TypingRule.SDecl.bad2.asl}


\RenderProseAndFormally{annotate_stmt_SDecl}
\CodeSubsection{\SDeclegin}{\SDeclEnd}{../Typing.ml}

\TypingRuleDef{AnnotateLocalDeclTypeAnnot}
\RenderRelation{annotate_local_decl_type_annot}

See \ExampleRef{Typing Declaration Statements}.


\RenderProseAndFormally{annotate_local_decl_type_annot}

\TypingRuleDef{InheritIntegerConstraints}
\RenderRelation{inherit_integer_constraints}

\listingref{typing-pendingconstrained} shows examples of pending-constrained
integer types.

\ExampleDef{Pending-constrained integer type vs. unconstrained integer type}
\listingref{pending-constrained-unconstrained} corresponds to the \typingerrorterm{}
in the \CaseName{int} below.
\ASLListing{An ill-typed pending-constrained integer type}{pending-constrained-unconstrained}
{\typingtests/TypingRule.InheritIntegerConstraints.unconstrained.bad.asl}


\RenderProseAndFormally{inherit_integer_constraints}

\TypingRuleDef{CheckNoPrecisionLoss}
\RenderRelation{check_no_precision_loss}

\ExampleDef{Rejected Declaration Because of Precision Loss}
In \listingref{typing-imprecisetype}, the statement \verb|var b = a * a;|
corresponds to the \typingerrorterm{} raised in the \CaseName{Well-Constrained}
below.
The type of the right-hand-side (\texttt{a * a}) is imprecise because the multiplication of two
constant types would result in more than $\maxconstraintsize$ elements, see
\TypingRuleRef{ApplyBinopTypes}.
\ASLListing{Typechecking a declaration with an imprecise type}{typing-imprecisetype}{\typingtests/TypingRule.CheckNoPrecisionLoss.asl}


\RenderProseAndFormally{check_no_precision_loss}
\CodeSubsection{\CheckNoPrecisionLossBegin}{\CheckNoPreciisonLossEnd}{../Typing.ml}

\TypingRuleDef{CheckCanBeInitializedWith}
\RenderRelation{check_can_be_initialized_with}

See \ExampleRef{Type-satisfaction Examples}.


\RenderProseAndFormally{check_can_be_initialized_with}


\subsection{Semantics}
\SemanticsRuleDef{SDecl}
Notice that \TypingRuleRef{SDecl} ensures that a typed \declarationstatementterm{}
has an initializing expression.
% The implementation still handles the case where there is no initializing expression,
% for cases where typechecking is skipped.

\ExampleDef{Declaration With an Initializing Value}
In \listingref{semantics-sdeclsome},
\texttt{let x = 3;} binds \texttt{x} to $\nvint(3)$ in the empty environment.
\ASLListing{Evaluating a declaration with a given initial value}{semantics-sdeclsome}{\semanticstests/SemanticsRule.SDeclSome.asl}

\ExampleDef{Declaration Without an Initializing Value}
In \listingref{semantics-sdeclnone},
\texttt{var x : integer;} binds \texttt{x} in $\env$ to the base value of \texttt{integer}.
\ASLListing{Evaluating a declaration without a given initial value}{semantics-sdeclnone}{\semanticstests/SemanticsRule.SDeclNone.asl}


\RenderProseAndFormally{eval_stmt_SDecl}
\CodeSubsection{\EvalSDeclBegin}{\EvalSDeclEnd}{../Interpreter.ml}

\section{Declaration Statements with an Elided Parameter\label{sec:DeclarationStatementsElidedParameter}}
\subsection{Syntax}
\begin{flalign*}
\Nstmt \derives \
   & \Nlocaldeclkeyword \parsesep \Ndeclitem \parsesep \Nasty \parsesep \Teq \\
   & \wrappedline\ \Nelidedparamcall \parsesep \Tsemicolon &\\
\end{flalign*}

\begin{flalign*}
\Nelidedparamcall \derives \
     & \Tidentifier \parsesep \Tlbrace \parsesep \Trbrace &\\
  |\ & \Tidentifier \parsesep \Tlbrace \parsesep \Trbrace \parsesep \PlistZero{\Nexpr} &\\
  |\ & \Tidentifier \parsesep \Tlbrace \parsesep \Tcomma \parsesep \ClistOne{\Nexpr} \parsesep \Trbrace &\\
  |\ & \Tidentifier \parsesep \Tlbrace \parsesep \Tcomma \parsesep \ClistOne{\Nexpr} \parsesep \Trbrace \parsesep \PlistZero{\Nexpr}&
\end{flalign*}

\SyntacticSugarDef{ParameterElision}
In a call to a parameterized subprogram, a parameter can be elided,
if it can be inferred from the left-hand-side type annotation,
as detailed by \ASTRuleRef{ElidedParamCall} (See also \secref{ParameterElision}).

\ASTRuleDef{ElidedParamCall}
\RenderRelation{build_elided_param_call}


\RenderProseAndFormally{build_elided_param_call}

\ASTRuleDef{DesugarElidedParameter}
\RenderRelation{desugar_elided_parameter}


\ExampleDef{Desugaring Parameter Elision Based on Declared Type Annotation}
\listingref{DesugarElidedParameter} shows examples of how parameters can be elided
if they can be copied from the type annotation in declaration statements.
\ASLListing{Desugaring parameter elision based on declared type annotation}{DesugarElidedParameter}{\syntaxtests/ASTRule.DesugarElidedParameter.asl}


\RenderProseAndFormally{desugar_elided_parameter}

\ASTRuleDef{ElidedParamDecl}

\RenderProseAndFormally{build_stmt_ElidedParamDecl}

\subsection{Typing and semantics}
As given by the applying the relevant rules to the desugared AST (see \secref{DeclarationStatements}).

\section{Sequencing Statements\label{sec:SequencingStatement}}
\hypertarget{def-sequencestatementterm}{}

\ASLListing{A sequence of statements}{semantics-sseq}{\semanticstests/SemanticsRule.SSeq.asl}

\subsection{Syntax}
\begin{flalign*}
\Nstmtlist \derives \ & \ListOne{\Nstmt} &
\end{flalign*}

\subsection{Abstract Syntax}
\RenderTypes[remove_hypertargets]{stmt_seq}


\ASTRuleDef{StmtList}
\RenderRelation{build_stmt_list}


\RenderProseAndFormally{build_stmt_list}

\ASTRuleDef{StmtFromList}
\RenderRelation{stmt_from_list}


\RenderProseAndFormally{stmt_from_list}

\ASTRuleDef{SequenceStmts}
\RenderRelation{sequence_stmts}


\RenderProseAndFormally{sequence_stmts}

\subsection{Typing}
\TypingRuleDef{SSeq}
\ExampleDef{Typing Sequencing Statements}
In \listingref{semantics-sseq}, the statement \verb|let x=3;| is
annotated first in the \staticenvironmentterm{} where the \localstaticenvironmentterm{} is empty.
Then, the statement \verb|let y = x + 1;|
is annotated in the \staticenvironmentterm{} where
\verb|x| has been declared and associated with the type \verb|integer{3}|,
and also recorded to be equivalent to $\LInt(3)$ (in the $\exprequiv$ map).


\RenderProseAndFormally{annotate_stmt_SSeq}
\CodeSubsection{\SSeqBegin}{\SSeqEnd}{../Typing.ml}

\subsection{Semantics}
\SemanticsRuleDef{SSeq}
\ExampleDef{Evaluation of Sequencing Statements}
In \listingref{semantics-sseq},
the evaluation of \texttt{let x = 3; let y = x + 1} first evaluates \texttt{let x = 3} and only then
evaluates \texttt{let y = x + 1}.


\RenderProseAndFormally{eval_stmt_SSeq}
\CodeSubsection{\EvalSSeqBegin}{\EvalSSeqEnd}{../Interpreter.ml}

\section{Call Statements\label{sec:CallStatements}}
\hypertarget{def-callstatementterm}{}
Call statements are used to invoke procedures and setters.

\ASLListing{Call Statements}{scall}{\semanticstests/SemanticsRule.SCall.asl}

\subsection{Syntax}
\begin{flalign*}
\Nstmt \derives \ & \Ncall \Tsemicolon &
\end{flalign*}

\subsection{Abstract Syntax}
\RenderTypes[remove_hypertargets]{stmt_call}

\ASTRuleDef{SCall}

\RenderProseAndFormally{build_stmt_SCall}

\ASTRuleDef{SetCallType}
\RenderRelation{set_call_type}
In \ASTRuleRef{Call}, $\STFunction$ is inserted as a default call type for any parsed $\call$.
The function
\[
  \setcalltype(\overname{\call}{\vcall} \aslsep \overname{\subprogramtype}{\calltype}) \aslto \overname{\call}{\vcallp}
\]
changes the call type of $\vcall$ to $\calltype$.


\RenderProseAndFormally{set_call_type}

\subsection{Typing}
\TypingRuleDef{SCall}
In \listingref{scall}, the call statement \verb|catenate_into_g{4, 3}(x, y, TRUE);|
is well-typed.

\ExampleDef{Ill-defined Call to a Function}
In \listingref{scall-bad}, the call statement \verb|zero();| is ill-typed,
since call statements are only allowed for procedures, not functions.

\ASLListing{An ill-defined call statement}{scall-bad}{\typingtests/TypingRule.SCall.bad.asl}


\RenderProseAndFormally{annotate_stmt_SCall}
\CodeSubsection{\SCallBegin}{\SCallEnd}{../Typing.ml}


\subsection{Semantics}
\SemanticsRuleDef{SCall}
\ExampleDef{Evaluation of Call Statements}
In \listingref{scall}, the call statement \verb|catenate_into_g{4, 3}(x, y, TRUE);|
assigns \verb|'1101 111'| to the global variable \verb|g|.


\RenderProseAndFormally{eval_stmt_SCall}
\CodeSubsection{\EvalSCallBegin}{\EvalSCallEnd}{../Interpreter.ml}


\hypertarget{def-conditionalstatementterm}{}
\section{Conditional Statements\label{sec:ConditionalStatements}}

A conditional statement evaluates its \texttt{then} statement list if the
condition expression evaluates to $\True$. If the condition expression
evaluates to $\False$ each \texttt{elsif} condition expression is evaluated
sequentially until an \texttt{elsif} condition expression evaluates to $\True$;
the conditional statement evaluates the corresponding \texttt{elsif}
statement list. If no \texttt{elsif} condition expression evaluates to $\True$, the
conditional statement evaluates the \texttt{else} statement list.

\subsection{Syntax}
\begin{flalign*}
\Nstmt \derives \ & \Tif \parsesep \Nexpr \parsesep \Tthen \parsesep \Nstmtlist \parsesep \Nselse \parsesep \Tend \parsesep \Tsemicolon &\\
\Nselse \derives\ & \Telseif \parsesep \Nexpr \parsesep \Tthen \parsesep \Nstmtlist \parsesep \Nselse &\\
        |\ & \Telse \parsesep \Nstmtlist &\\
        |\ & \emptysentence &
\end{flalign*}

\subsection{Abstract Syntax}
\RenderTypes[remove_hypertargets]{stmt_cond}


\ASTRuleDef{SCond}

\RenderProseAndFormally{build_stmt_SCond}

\ASTRuleDef{SElse}
\RenderRelation{build_s_else}


\RenderProseAndFormally{build_s_else}

\subsection{Typing}
\TypingRuleDef{SCond}
The specifications in \listingref{semantics-scond},
\listingref{semantics-scond2},
\listingref{semantics-scond3}, and
\listingref{semantics-scond4} are all well-typed.

\ExampleDef{Typing Conditional Statements with Two Bitvector Types}
The specification in \listingref{SCond} is well-typed
and shows examples of how \bitvectortypesterm{} either \typesatisfyterm{}
the return type or do not \typesatisfyterm{} the return type and require
an \atcexpressionterm{}.
\ASLListing{Typing conditional statements with two bitvector types}{SCond}{\typingtests/TypingRule.SCond.asl}


\RenderProseAndFormally{annotate_stmt_SCond}
\CodeSubsection{\SCondBegin}{\SCondEnd}{../Typing.ml}


\subsection{Semantics}
\SemanticsRuleDef{SCond}
\ExampleDef{Conditional Statements}
The specification in \listingref{semantics-scond} does not result in any Assertion Error.
\ASLListing{Evaluating a conditional statement}{semantics-scond}{\semanticstests/SemanticsRule.SCond.asl}

The specification in \listingref{semantics-scond2}
does not result in any error.
\ASLListing{Evaluating a condition statement with \texttt{elsif}}{semantics-scond2}{\semanticstests/SemanticsRule.SCond2.asl}

The specification in \listingref{semantics-scond3}
results in an Assertion Error.
\ASLListing{Evaluating a condition statement that results in an Assertion Error}{semantics-scond3}{\semanticstests/SemanticsRule.SCond3.asl}

The specification in \listingref{semantics-scond4} does not result in any error.
\ASLListing{Evaluating a condition statement with only a \texttt{then} branch}{semantics-scond4}{\semanticstests/SemanticsRule.SCond4.asl}


\RenderProseAndFormally{eval_stmt_SCond}
\CodeSubsection{\EvalSCondBegin}{\EvalSCondEnd}{../Interpreter.ml}

\hypertarget{def-casestatementterm}{}
\section{Case Statements\label{sec:CaseStatements}}
Case statements allow executing different statements, based on which
condition an expression satisfies.

\listingref{case-discriminant} shows an example of a \casestatementterm{}
and the output to a console when it is evaluated.

\ASLListing{A side-effecting case discriminant}{case-discriminant}{\definitiontests/CaseStatement.discriminant.asl}
% CONSOLE_BEGIN aslref \definitiontests/CaseStatement.discriminant.asl
\begin{Verbatim}[fontsize=\footnotesize, frame=single]
num_tests: 0
selected case x=51
\end{Verbatim}
% CONSOLE_END
\listingref{case-discriminant} shows an example of a \casestatementterm{}
and the output to a console when it is evaluated.
%
\listingref{CaseStatement.where} shows an example of a \casestatementterm{}
with a \texttt{where} guard expression.

\hypertarget{def-casediscriminantterm}{}
\hypertarget{def-casealternativeterm}{}
\hypertarget{def-otherwisecaseterm}{}
The expression following the \Tcase\ keyword is called the \emph{\casediscriminantterm}.
The list following the \Tof\ keyword consists of \emph{\casealternativesterm},
optionally ending with an \emph{\otherwisecaseterm}, which follows the \Totherwise\ keyword.

Case statements obey the following requirements:

\RequirementDef{CaseDiscriminant} The \casediscriminantterm\ of a \texttt{case}
statement should be evaluated only once each time the case statement is evaluated.
%
\listingref{case-discriminant} demonstrates how the \casediscriminantterm{}
is evaluated only once.

\RequirementDef{CaseAlternatives} The \casealternativesterm\ are examined
one after another, in the order they are listed.
If any of the patterns match the \casediscriminantterm{} (and the guard
expression is true, if present) then this \casealternativeterm{} is considered selected,
its statement list is executed, and the \texttt{case} statement ends without examining any further
\casealternativesterm.
%
\listingref{case-discriminant} demonstrates how only one \casealternativeterm{} is selected
for execution.

\RequirementDef{CaseDiscriminantTesting} Testing the \casediscriminantterm{} against a \\
pattern list
follows the semantics of pattern matching defined in \chapref{PatternMatching}.
It is not a static error if it can be statically determined that none of the patterns in a
\casealternativeterm\ can match the discriminant.
%
\listingref{case-discriminant} exemplifies a \casestatementterm{} with pattern matching.

\RequirementDef{CaseOtherwise} If no \casealternativeterm\ is selected, and there is an
\otherwisecaseterm\, the \otherwisecaseterm\ is executed.
%
\listingref{case-otherwise} demonstrates how the \otherwisecaseterm{} is evaluated.

\ASLListing{Executing the \texttt{otherwise} case alternative}{case-otherwise}{\definitiontests/CaseStatement.otherwise.asl}
% CONSOLE_BEGIN aslref \definitiontests/CaseStatement.otherwise.asl
\begin{Verbatim}[fontsize=\footnotesize, frame=single]
num_tests: 0
selected otherwise
\end{Verbatim}
% CONSOLE_END

\RequirementDef{CaseNoOtherwiseError} If no \casealternativeterm\ is selected,
and there is no \otherwisecaseterm, it is a \dynamicerrorterm{}.
%
\listingref{case-no-otherwise} shows a specification that, when evaluated,
yields a \dynamicerrorterm{}.

\ASLListing{A case statement with no \texttt{otherwise} case}{case-no-otherwise}{\definitiontests/CaseStatement.no_otherwise.asl}

\subsection{Syntax}
\begin{flalign*}
\Nstmt \derives \ & \Tcase \parsesep \Nexpr \parsesep \Tof \parsesep \Ncasealtlist \parsesep \Tend \parsesep \Tsemicolon &\\
|\ & \Tcase \parsesep \Nexpr \parsesep \Tof \parsesep \Ncasealtlist \parsesep \Totherwise \parsesep \Tarrow &\\
   & \wrappedline\ \Nstmtlist \parsesep \Tend \parsesep \Tsemicolon &\\
\Ncasealtlist \derives \ & \ClistOne{\Ncasealt} \parsesep &\\
\Ncasealt \derives \ & \Twhen \parsesep \Npatternlist \parsesep \option{\Twhere \parsesep \Nexpr} \parsesep \Tarrow \parsesep \Nstmtlist &\\
\end{flalign*}

\listingref{CaseStatement-bad} shows an example of an erroneous case statement
where a \casealternativeterm{} is missing a statement.
\ASLListing{Invalid case statement syntax}{CaseStatement-bad}{\syntaxtests/CaseStatement.bad.asl}

\subsection{Abstract Syntax}

\SyntacticSugarDef{CaseStatement}
Case statements are considered \\ \syntacticsugar{}.
That is, there is no explicit representation for \texttt{case} statements
as they are represented in AST via a \conditionalstatementterm.
Case statements are \desugared{} by \ASTRuleRef{DesugarCaseStmt}.

\ExampleDef{Case Statement Desugaring}
\listingref{ast-case1} shows an example of how a \texttt{case} statement can be transformed into a corresponding
compound condition statement.
\ASLListing{Transforming a case statement with a variable as the case discriminant}{ast-case1}{\syntaxtests/ASTRule.Desugar_SCase1.asl}

\listingref{ast-case2} shows an example of how a \texttt{case} statement can be transformed into
a statement that does not contain any \text{case} statement.
By storing the \casediscriminantterm\ in a variable and by adding
an \unreachablestatementterm{}, the transformation ensures that a \dynamicerrorterm{} occurs when no
\casealternativeterm\ is selected.
\ASLListing{Transforming a case statement with a non-variable as the case discriminant
and no otherwise case}{ast-case2}{\syntaxtests/ASTRule.Desugar_SCase2.asl}

The untyped AST contains non-terminals for \casealternativesterm, which exist
only as a data type used by $\desugarcasestmt$ and do not later appear in the untyped
AST:

\RenderType[remove_hypertargets]{case_alt}

\ASTRuleDef{SCase}
To satisfy \RequirementRef{CaseNoOtherwiseError}, when no \otherwisecaseterm\ exists,\linebreak[2]
$\SUnreachable$ is used instead.
\RenderProseAndFormally{build_stmt_SCase}

\ASTRuleDef{CaseAltList}
\RenderRelation{build_case_alt_list}


\RenderProseAndFormally{build_case_alt_list}

\ASTRuleDef{CaseAlt}
\RenderRelation{build_case_alt}


\RenderProseAndFormally{build_case_alt}

\ASTRuleDef{DesugarCaseStmt}
\RenderRelation{desugar_case_stmt}

To satisfy \RequirementRef{CaseDiscriminant}, the transformation assigns the
\casediscriminantterm{} to a temporary variable, which is then used in a
compound conditional statement (see \listingref{ast-case2} for an example):


\RenderProseAndFormally{desugar_case_stmt}
\CodeSubsection{\DesugarCaseStmtBegin}{\DesugarCaseStmtEnd}{../desugar.ml}

\ASTRuleDef{CasesToCond}
\RenderRelation{cases_to_cond}


\RenderProseAndFormally{cases_to_cond}
\CodeSubsection{\CasesToCondBegin}{\CasesToCondEnd}{../desugar.ml}

\ASTRuleDef{CaseToCond}
\RenderRelation{case_to_cond}


\RenderProseAndFormally{case_to_cond}
\CodeSubsection{\CaseToCondBegin}{\CaseToCondEnd}{../desugar.ml}

\subsection{Typing}
Since case statements are transformed into other statements,
they do not require type system rules.


\subsection{Semantics}
Since case statements are transformed into other statements,
they do not appear in the typed AST and thus are not associated with a semantics.

In \listingref{CaseStatement.where} the expression \verb|a| will be evaluated twice.
\ASLListing{Evaluating case alternative sub-expressions}{CaseStatement.where}{\definitiontests/CaseStatement.where.asl}

\hypertarget{def-assertionstatementterm}{}
\section{Assertion Statements\label{sec:AssertionStatements}}
Assertion statements are used to check that certain conditions are satisfied.
They take a single \booleantypeterm{} operand, which we refer to as the
\emph{condition}. If the condition is \False, the statement fails with a
\dynamicerrorterm{} (\errorcodeterm{} $\DynamicAssertionFailure$).

\listingref{AssertionStatement} shows a possible use of
\assertionstatementterm{} to check the inputs to a function
(also known as a \emph{precondition})
and ensure the output satisfies expectations (also known as a \emph{postcondition}).
%
The first call to \verb|checked_8bit_add| succeeds, whereas the second call fails the
\assertionstatementterm{} \verb|assert a + b < 256;|

\ASLListing{Example of using \texttt{assert} statements}{AssertionStatement}{\definitiontests/AssertionStatement.asl}

\subsection{Syntax}
\begin{flalign*}
\Nstmt \derives \ & \Tassert \parsesep \Nexpr \parsesep \Tsemicolon &
\end{flalign*}

\subsection{Abstract Syntax}
\RenderTypes[remove_hypertargets]{stmt_assert}


\ASTRuleDef{SAssert}

\RenderProseAndFormally{build_stmt_SAssert}

\subsection{Typing}
\TypingRuleDef{SAssert}
\ExampleDef{Type Assertion Statements}
The specifications in \listingref{semantics-sasertok} and \listingref{semantics-sassertno}
are well-typed.

The assertion statements in \listingref{assert-ill-typed} are ill-typed,
since the expression needs to be both of a \booleantypeterm{} and \readonlyterm{}.
\ASLListing{Ill-typed assertion statements}{assert-ill-typed}{\typingtests/TypingRule.SAssert.bad.asl}


\RenderProseAndFormally{annotate_stmt_SAssert}
\CodeSubsection{\SAssertBegin}{\SAssertEnd}{../Typing.ml}


\subsection{Semantics}
\SemanticsRuleDef{SAssert}
\ExampleDef{Evaluation of Assertions: Success}
In \listingref{semantics-sasertok},
\texttt{assert (42 != 3);} ensures that \texttt{3} is not equal to \texttt{42}.
\ASLListing{Evaluating an assertion that succeeds}{semantics-sasertok}{\semanticstests/SemanticsRule.SAssertOk.asl}

\ExampleDef{Evaluation of Assertions: Failure}
In \listingref{semantics-sassertno},
evaluating \texttt{assert (42 == 3);} results in an \DynamicAssertionFailure{} error.
\ASLListing{Evaluating an assertion that fails}{semantics-sassertno}{\semanticstests/SemanticsRule.SAssertNo.asl}


\RenderProseAndFormally{eval_stmt_SAssert}
\CodeSubsection{\EvalSAssertBegin}{\EvalSAssertEnd}{../Interpreter.ml}


\section{While Statements\label{sec:WhileStatements}}
\hypertarget{def-whilestatementterm}{}

\ASLListing{A \texttt{while} statement}{semantics-swhile}{\semanticstests/SemanticsRule.SWhile.asl}

\ConventionDef{Loop Limits}
Conventionally, all loop kinds should specify a limit expression.
For example, the loops in \listingref{semantics-swhile} specifies a limit via
the expression \verb|limit_loop()|.

\subsection{Syntax}
\begin{flalign*}
\Nstmt \derives \ & \Twhile \parsesep \Nexpr \parsesep \Nlooplimit \parsesep \Tdo \parsesep \Nstmtlist \parsesep \Tend \parsesep \Tsemicolon &
\end{flalign*}

\subsection{Abstract Syntax}
\RenderTypes[remove_hypertargets]{stmt_while}


\ASTRuleDef{SWhile}

\RenderProseAndFormally{build_stmt_SWhile}

\ASTRuleDef{LoopLimit}
\RenderRelation{build_loop_limit}


\RenderProseAndFormally{build_loop_limit}

\subsection{Typing}
\TypingRuleDef{SWhile}
\ExampleDef{Typing While Loops}
The specification in \listingref{typing-swhile} is well-typed
and shows two \whilestatementterm{} --- the one in \verb|scan| contains a loop limit
and the one in \verb|main| does not contain a loop limit.
\ASLListing{Typing while statements}{typing-swhile}{\typingtests/TypingRule.SWhile.asl}

The specification in \listingref{typing-swhile-bad-limit}
is ill-typed, since the loop limit is not a \constrainedintegerterm.
\ASLListing{An ill-typed loop limit}{typing-swhile-bad-limit}{\typingtests/TypingRule.SWhile.bad_limit.asl}


\RenderProseAndFormally{annotate_stmt_SWhile}
\CodeSubsection{\SWhileBegin}{\SWhileEnd}{../Typing.ml}


\TypingRuleDef{AnnotateLimitExpr}
\RenderRelation{annotate_limit_expr}

\ExampleDef{Annotating Limit Expressions}
The specification in \listingref{typing-swhile} has two loops.
The loop in \verb|scan| is limited by the constrained integer
expression \verb|N|, while the loop in \verb|main| does not
have a loop limit.

The loop in \listingref{typing-swhile-bad-limit} uses
the limit expression \verb|i_limit| whose type the
\unconstrainedintegertypeterm, which is illegal for
limit expressions.


\RenderProseAndFormally{annotate_limit_expr}
\CodeSubsection{\AnnotateLimitExprBegin}{\AnnotateLimitExprEnd}{../Typing.ml}

\subsection{Semantics}
\SemanticsRuleDef{SWhile}
\ExampleDef{Evaluation of While Statements}
The specification in \listingref{semantics-swhile} prints the
following to the console:
% CONSOLE_BEGIN aslref \semanticstests/SemanticsRule.SWhile.asl
\begin{Verbatim}[fontsize=\footnotesize, frame=single]
i = 0
i = 1
i = 2
i = 3
\end{Verbatim}
% CONSOLE_END

The specification in \listingref{swhile-limit-reached} yields a dynamic
error once the loop limit for the \whilestatementterm{} is reached.
\ASLListing{Reaching a while loop limit}{swhile-limit-reached}{\semanticstests/SemanticsRule.SWhile.limit_reached.asl}


\RenderProseAndFormally{eval_stmt_SWhile}
\CodeSubsection{\EvalSWhileBegin}{\EvalSWhileEnd}{../Interpreter.ml}

\SemanticsRuleDef{Loop}
\RenderRelation{eval_loop}

More specifically, $\evalloop(\env, \iswhile, \velimitopt, \econd, \vbody)$
evaluates\\ $\vbody$ in $\env$ as long as $\econd$ holds when $\iswhile$ is $\True$
or until $\econd$ holds when $\iswhile$ is $\False$.
If the number of iterations exceeds the optional value specified by $\vlimitopt$,
the result is a \dynamicerrorterm{}.
The result is either the continuing configuration \\ $\Continuing(\newg,\newenv)$,
an early return configuration, or an abnormal configuration.

\ExampleDef{Evaluation of Loops}
The specification in \listingref{semantics-loop} does not result in any Assertion Error
and the specification terminates with exit code $0$.
\ASLListing{Evaluating a loop}{semantics-loop}{\semanticstests/SemanticsRule.Loop.asl}

In the following definition,
the premise $\vb \neq \iswhile$ is $\True$ in the case of a \texttt{while} loop
and the loop condition $\ve$ not holding, which is exactly when we want the
loop to exit. The opposite holds for a \texttt{repeat} loop.
The negation of the condition is used to decide whether to continue the loop iteration.

\RenderProseAndFormally{eval_loop}
\CodeSubsection{\EvalLoopBegin}{\EvalLoopEnd}{../Interpreter.ml}

\SemanticsRuleDef{EvalLimit}
\RenderRelation{eval_limit}

The evaluation uses the function $\evalexprsef$ because limit expressions are
guaranteed \sideeffectfreeterm{} by the typechecker,
see \TypingRuleRef{AnnotateLimitExpr}.

See \ExampleRef{Evaluation of While Statements} to see how the
limit expression is evaluated just once for the entire evaluation
of the \whilestatementterm.


\RenderProseAndFormally{eval_limit}

\SemanticsRuleDef{TickLoopLimit}
\RenderRelation{tick_loop_limit}

\ExampleDef{Decrementing a Loop Limit Value}
In \listingref{semantics-swhile}, the loop limit value
starts at $\nvint(4)$ and decrements towards $\nvint(0)$,
stopping at $\nvint(1)$ on the last iteration of the loop.

\ExampleDef{Negative Loop Limit Values}
In \listingref{semantics-swhile-negative-limit}, the loop limit value starts at $\nvint(-1)$
so the loop limit is immediately reached, resulting in a \dynamicerrorterm{} without executing the loop body at all.
(The same would happen if the loop limit value started at $\nvint(0)$.)
\ASLListing{A \texttt{while} statement with a negative loop limit}{semantics-swhile-negative-limit}{\semanticstests/SemanticsRule.SWhile.negative_limit.asl}

\RenderProseAndFormally{tick_loop_limit}

\section{Repeat Statements\label{sec:RepeatStatements}}
\hypertarget{def-repeatstatementterm}{}

\ASLListing{A repeat statement}{semantics-srepeat}{\semanticstests/SemanticsRule.SRepeat.asl}

\subsection{Syntax}
\begin{flalign*}
\Nstmt \derives \ & \Trepeat \parsesep \Nstmtlist \parsesep \Tuntil \parsesep \Nexpr \parsesep \Nlooplimit \parsesep \Tsemicolon &
\end{flalign*}

\subsection{Abstract Syntax}
\RenderTypes[remove_hypertargets]{stmt_repeat}

\ASTRuleDef{SRepeat}

\RenderProseAndFormally{build_stmt_SRepeat}

\subsection{Typing}
\TypingRuleDef{SRepeat}
\ExampleDef{Typing a Repeat Statement}
The \repeatstatementsterm{} in \listingref{semantics-srepeat} are well-typed.


\RenderProseAndFormally{annotate_stmt_SRepeat}
\CodeSubsection{\SRepeatBegin}{\SRepeatEnd}{../Typing.ml}


\subsection{Semantics}
\SemanticsRuleDef{SRepeat}
\ExampleDef{Evaluation of Repeat Statements}
The specification in \listingref{semantics-srepeat} produces the following output to the console.
% CONSOLE_BEGIN aslref  --type-check-no-warn \semanticstests/SemanticsRule.SRepeat.asl
\begin{Verbatim}[fontsize=\footnotesize, frame=single]
j = 0
j = 1
j = 2
j = 3
j = 4
#ones in x = 5
i = 0
i = 1
i = 2
i = 3
i = 4
#ones in x = 5
\end{Verbatim}
% CONSOLE_END


\RenderProseAndFormally{eval_stmt_SRepeat}
\CodeSubsection{\EvalSRepeatBegin}{\EvalSRepeatEnd}{../Interpreter.ml}

\section{For Statements\label{sec:ForStatements}}
\hypertarget{def-forstatementterm}{}

\ASLListing{\texttt{for} loops}{semantics-sfor}{\semanticstests/SemanticsRule.SFor.asl}

\subsection{Syntax}
\begin{flalign*}
\Nstmt \derives \ & \Tfor \parsesep \Tidentifier \parsesep \Teq \parsesep \Nexpr \parsesep \Ndirection \parsesep
                    \Nexpr \parsesep \Nlooplimit \parsesep \Tdo &\\
                        & \wrappedline\ \Nstmtlist \parsesep \Tend \parsesep \Tsemicolon &\\
\Ndirection \derives \ & \Tto \;|\; \Tdownto &
\end{flalign*}

\subsection{Abstract Syntax}
\RenderTypes[remove_hypertargets]{stmt_for}


\ASTRuleDef{SFor}

\RenderProseAndFormally{build_stmt_SFor}

\ASTRuleDef{Direction}
\RenderRelation{build_direction}


\RenderProseAndFormally{build_direction}

\subsection{Typing}
\TypingRuleDef{SFor}
The \texttt{for} loops in \listingref{semantics-sfor} are well-typed.

\ExampleDef{Ill-typed for loop}
\listingref{typing-forloop-bad1},
\listingref{typing-forloop-bad2},
\listingref{typing-forloop-bad3}, and
\listingref{typing-forloop-bad4}
show examples of ill-typed \texttt{for} loops.

\ASLListing{Ill-typed for loop 1}{typing-forloop-bad1}{\typingtests/TypingRule.SFor.bad1.asl}
\ASLListing{Ill-typed for loop 2}{typing-forloop-bad2}{\typingtests/TypingRule.SFor.bad2.asl}
\ASLListing{Ill-typed for loop 3}{typing-forloop-bad3}{\typingtests/TypingRule.SFor.bad3.asl}
\ASLListing{Ill-typed for loop 4}{typing-forloop-bad4}{\typingtests/TypingRule.SFor.bad4.asl}


\RenderProseAndFormally{annotate_stmt_SFor}
\CodeSubsection{\SForBegin}{\SForEnd}{../Typing.ml}


\TypingRuleDef{SForConstraints}
\RenderRelation{get_for_constraints}

\ExampleDef{Inferring the Constraints of a for Loop Index}
In \listingref{semantics-sfor} the constraints for the loop index variable
\verb|j| in \verb|scan| are \verb|0..N-1|,
and the constraints for the loop index variable \verb|i| in \verb|main| are \verb|0..4|.


\RenderProseAndFormally{get_for_constraints}

\subsection{Semantics}
\SemanticsRuleDef{SFor}
Evaluating a \texttt{for} statement involves introducing an index variable to the
environment. The type system ensures, via \TypingRuleRef{SFor}, that the index variable
is not already declared in the scope of the subprogram containing the \texttt{for}
statement.

\pagebreak % To avoid cutting the example output.
\ExampleDef{Evaluation of For Statements}
The specification in \listingref{semantics-sfor} is followed by its output to the console.
% CONSOLE_BEGIN aslref \semanticstests/SemanticsRule.SFor.asl
\begin{Verbatim}[fontsize=\footnotesize, frame=single]
j = 0
j = 1
j = 2
j = 3
j = 4
#ones in x = 5
i = 4
i = 3
i = 2
i = 1
i = 0
#ones in x = 5
\end{Verbatim}
% CONSOLE_END

The specification in \listingref{SFor-nop} shows a \forstatementterm{}
whose body does not evaluate since its lower bound is greater than its upper bound.
\ASLListing{A for loop with no iterations}{SFor-nop}{\semanticstests/SemanticsRule.SFor.nop.asl}

Recall that the expressions for the \texttt{for} loop range are
\sideeffectfreeterm{}, as guaranteed by \TypingRuleRef{SFor}, which is why
they are evaluated via $\evalexprsef$.

\RenderProseAndFormally{eval_stmt_SFor}
\CodeSubsection{\EvalSForBegin}{\EvalSForEnd}{../Interpreter.ml}

\SemanticsRuleDef{EvalFor}
\RenderRelation{eval_for}

\ExampleDef{Overall Evaluation of For Statements}
The specification in \listingref{semantics-sfor}
does not result in any assertion error, and terminates with exit-code $0$.

Evaluating $(\vindexname, \vstart, \dir, \vend, \vbody)$ in $\env$ yields either
a continuing configuration $\Continuing(\newg, \newenv)$, or a returning configuration
(in case the body of the loop results in an early return),
or an abnormal configuration.


\RenderProseAndFormally{eval_for}
\CodeSubsection{\EvalForBegin}{\EvalForEnd}{../Interpreter.ml}

\SemanticsRuleDef{EvalForLoop}
\RenderRelation{eval_for_loop}

See \ExampleRef{Overall Evaluation of For Statements}.


\RenderProseAndFormally{eval_for_loop}

\SemanticsRuleDef{EvalForStep}
\RenderRelation{eval_for_step}

See \ExampleRef{Overall Evaluation of For Statements}.

\RenderProseAndFormally{eval_for_step}

\hypertarget{def-throwstatementterm}{}
\section{Throw Statements\label{sec:ThrowStatements}}
\subsection{Syntax}
\begin{flalign*}
\Nstmt \derives \ & \Tthrow \parsesep \Nexpr \parsesep \Tsemicolon &
\end{flalign*}

\subsection{Abstract Syntax}
\RenderTypes[remove_hypertargets]{stmt_throw}


\ASTRuleDef{SThrow}

\RenderProseAndFormally{build_stmt_SThrow}

\subsection{Typing}
\TypingRuleDef{SThrow}
\listingref{semantics-sthrowsometyped}
shows examples of well-typed \throwstatementsterm.


\RenderProseAndFormally{annotate_stmt_SThrow}

\CodeSubsection{\SThrowBegin}{\SThrowEnd}{../Typing.ml}

\subsection{Semantics}
\SemanticsRuleDef{SThrow}
\ExampleDef{Throwing a Typed Exception}
The specification in \listingref{semantics-sthrowsometyped}
terminates successfully. That is, no \dynamicerrorterm{} occurs.
\ASLListing{Throwing an exception}{semantics-sthrowsometyped}{\semanticstests/SemanticsRule.SThrowSomeTyped.asl}


\RenderProseAndFormally{eval_stmt_SThrow}
\CodeSubsection{\EvalSThrowBegin}{\EvalSThrowEnd}{../Interpreter.ml}

\hypertarget{def-trystatementterm}{}
\section{Try Statements\label{sec:TryStatements}}
\subsection{Syntax}
\begin{flalign*}
\Nstmt \derives \ & \Ttry \parsesep \Nstmtlist \parsesep \Tcatch \parsesep \ListOne{\Ncatcher} \parsesep \Notherwiseopt &\\
                  & \wrappedline\ \parsesep \Tend \parsesep \Tsemicolon &\\
\Notherwiseopt \derives\ & \Totherwise \parsesep \Tarrow \parsesep \Nstmtlist &\\
               |\ & \emptysentence &\\
\end{flalign*}

\subsection{Abstract Syntax}
\RenderTypes[remove_hypertargets]{stmt_try}


\ASTRuleDef{STry}

\RenderProseAndFormally{build_stmt_STry}

\ASTRuleDef{OtherwiseOpt}
\RenderRelation{build_otherwise_opt}


\RenderProseAndFormally{build_otherwise_opt}

\subsection{Typing}

\TypingRuleDef{STry}
\ExampleDef{Typing Try Statements}
In
\listingref{semantics-catchnamed},
\listingref{semantics-catchotherwise},
\listingref{semantics-catchnone},
\listingref{semantics-nothrow}, and
\listingref{semantics-stry},
the \trystatementsterm{} are all well-typed.

\RenderProseAndFormally{annotate_stmt_STry}
\CodeSubsection{\STryBegin}{\STryEnd}{../Typing.ml}

\subsection{Semantics}
\SemanticsRuleDef{STry}
\ExampleDef{Evaluation of Try Statements}
Evaluating the specification in \listingref{semantics-stry}
does not result in any Assertion error, and the specification terminates with the exit code $0$.
\ASLListing{Evaluating a \texttt{try} statement}{semantics-stry}{\semanticstests/SemanticsRule.STry.asl}


\RenderProseAndFormally{eval_stmt_STry}
\CodeSubsection{\EvalSTryBegin}{\EvalSTryEnd}{../Interpreter.ml}

\hypertarget{def-returnstatementterm}{}
\section{Return Statements\label{sec:ReturnStatements}}
\subsection{Syntax}
\begin{flalign*}
\Nstmt \derives \ & \Treturn \parsesep \option{\Nexpr} \parsesep \Tsemicolon &
\end{flalign*}

\subsection{Abstract Syntax}
\RenderTypes[remove_hypertargets]{stmt_return}

\ASTRuleDef{SReturn}

\RenderProseAndFormally{build_stmt_SReturn}

\subsection{Typing}
\TypingRuleDef{SReturn}

\ExampleDef{Typing Return Statements}
The \returnstatementsterm{} in
\listingref{semantics-sreturn},
\listingref{semantic-ssreturnone}, and
\listingref{semantics-sreturntuple}
are all well-typed.

The return statement \verb|return 0;| in \listingref{typing-return-bad} is ill-typed,
since \verb|proc| is not a function but a procedure.
\ASLListing{An ill-typed return statement}{typing-return-bad}{\typingtests/TypingRule.SReturn.bad.asl}


\RenderProseAndFormally{annotate_stmt_SReturn}
\CodeSubsection{\SReturn}{\SReturnEnd}{../Typing.ml}


\subsection{Semantics}
\SemanticsRuleDef{SReturn}
\ExampleDef{No Return Value}
The specification in \listingref{semantics-sreturn} exits the procedure \texttt{println\_me}
by evaluating the\\ \texttt{return;} statement.
\ASLListing{Evaluating a \texttt{return} statement with no value}{semantics-sreturn}{\semanticstests/SemanticsRule.SReturnNone.asl}

\ExampleDef{Returning a Single Value}
In \listingref{semantic-ssreturnone},
\texttt{return 3;} exits the function \texttt{f} with value \texttt{3}.
\ASLListing{Evaluating a \texttt{return} statement with a single value}{semantic-ssreturnone}{\semanticstests/SemanticsRule.SReturnOne.asl}

\ExampleDef{Returning a Tuple of Values}
In \listingref{semantics-sreturntuple},
\texttt{return (3, 42);} exits the function \texttt{f} with the value \texttt{(3, 42)}.
\ASLListing{Evaluating a \texttt{return} statement with a tuple of values}{semantics-sreturntuple}{\semanticstests/SemanticsRule.SReturnSome.asl}


\RenderProseAndFormally{eval_stmt_SReturn}
\CodeSubsection{\SReturnBegin}{\EvalSReturnEnd}{../Interpreter.ml}

\SemanticsRuleDef{EExprListM}
\RenderRelation{eval_expr_list_m}

\ExampleDef{Seperately Evaluating a List of Expressions}
In \listingref{semantics-sreturntuple}, the expressions \verb|3| and \verb|42|
are evaluated in left-to-right order in the statement \verb|return (3 , 42);|.


\RenderProseAndFormally{eval_expr_list_m}

\SemanticsRuleDef{WriteFolder}
\RenderRelation{write_folder}

\ExampleDef{Folding a List of Pairs with Values and Execution Graphs}
In \listingref{semantics-sreturntuple}, the statement \verb|return (3 , 42);|
uses $\writefolder$ to generate \\
$([\nvint(3), \nvint(42)], \emptygraph)$.
The \executiongraphterm{} is empty, since literal expressions do not yield
\executiongraphsterm{} and composing empty \executiongraphsterm{} yields an
empty \executiongraphterm{}.


\RenderProseAndFormally{write_folder}

\hypertarget{def-printstatementterm}{}
\section{Print Statements\label{sec:PrintStatements}}
\subsection{Syntax}
\begin{flalign*}
\Nstmt \derives \ & \Tprint \parsesep \ClistZero{\Nexpr} \parsesep \Tsemicolon & \\
\Nstmt \derives \ & \Tprintln \parsesep \ClistZero{\Nexpr} \parsesep \Tsemicolon & \\
\end{flalign*}

\subsection{Abstract Syntax}
\RenderTypes[remove_hypertargets]{stmt_print}


\ASTRuleDef{SPrint}

\RenderProseAndFormally{build_stmt_SPrint}

\subsection{Typing}
\TypingRuleDef{SPrint}
\listingref{literals1} shows literals and their corresponding types in comments.


\RenderProseAndFormally{annotate_stmt_SPrint}
\CodeSubsection{\SPrintBegin}{\SPrintEnd}{../Typing.ml}

\subsection{Semantics}
\SemanticsRuleDef{SPrint}

\ExampleDef{Printing Literals}
\listingref{semantics-literals} shows examples of printing various types of literals,
followed by the output to the console resulting from running the specification.
\ASLListing{Literals and how they are displayed}{semantics-literals}{\semanticstests/SemanticsRule.SPrint.asl}
% CONSOLE_BEGIN aslref \semanticstests/SemanticsRule.SPrint.asl
\begin{Verbatim}[fontsize=\footnotesize, frame=single]
string_number_1
0
1000000
53170898287292728730499578000
TRUE
FALSE
12345678900123456789/10000000000
0
hello\world
	 "here I am "
0xd
0x
LABEL_B
\end{Verbatim}
% CONSOLE_END

Notice that empty bitvectors are displayed as \verb|0x|.

We define the constant \RenderConstant{new_line}, which stands
for the newline character: $\vnewline \triangleq \ascii{10}$.

\RenderProseAndFormally{eval_stmt_SPrint}
\CodeSubsection{\EvalSPrintBegin}{\EvalSPrintEnd}{../Interpreter.ml}

Not all ASL runtimes support printing to a console (see \RequirementRef{Printing}).
%
Therefore, the semantics are parameterized by the following function.
\RenderRelation{output_to_console}

We now explain how printing is modelled when the runtime supports a console
(\SemanticsRuleRef{SupportedOutputToConsole})
and how it is modelled when the runtime does not support a console
(\SemanticsRuleRef{UnsupportedOutputToConsole}).

\SemanticsRuleDef{SupportedOutputToConsole}
To support a console, the definition of environments needs
to include an extra component to capture the string of characters sent to the console:
\[
\envs \triangleq \staticenvs \cartimes \dynamicenvs \cartimes \Strings \enspace.
\]
We omit this component in the rest of this document to avoid clutter, and include it
only here to explain the modeling of a console.

\ExampleRef{Printing Literals} shows the output to the console in a case it is supported.

\RenderRelation{literal_to_string}
The function is defined by \taref{literaltostringtable}.
%
Please note that surrounding quotations marks that appear in ASL text (for example, \verb|"hello"|)
are not included in string literals, and thus will not appear in the printed string.

\begin{table}
\caption{How literals should be represented as strings\label{ta:literaltostringtable}}
\begin{tabular}{rl}
\textbf{literal $\vl$} & \textbf{$\literaltostring(\vl)$} \\
\hline
$\LInt(n)$        & $n$ in decimal format, without any leading zeros, \\
                  & preceded by a ``\texttt{-}'' sign if $n$ is negative. \\
$\LBool(\True)$   & \texttt{TRUE} \\
$\LBool(\False)$  & \texttt{FALSE} \\
$\LReal(q)$       & $q$ as an irreducible fraction of positive integers, \\
                  & preceded by a ``\texttt{-}'' sign when $q$ is negative, \\
                  & with the denominator omitted if it is equal to 1. \\
$\LBitvector(b)$  & $b$ in hexadecimal, preceded by ``\texttt{0x}'', with enough leading \\
                  & zeros to make the number of hexadecimal digits printed \\
                  & equal to the width of $b$ divided by 4, and rounded up to \\
                  & the following integer.\\
$\LString(S)$     & $S$. \\
$\LLabel(s)$      & $s$. \\
\end{tabular}
\end{table}

\ProseParagraph
\AllApply
\begin{itemize}
  \item view $\env$ as the environment consisting of the \staticenvironmentterm{} $\tenv$,
        dynamic environment $\denv$, and console string $\vconsolestream$;
  \item view $\vv$ as a native literal for the literal $\vl$;
  \item \Proseeqdef{$\newenv$}{the environment consisting of the \staticenvironmentterm{} $\tenv$,
        dynamic environment $\denv$, and console string
        define as the concatenation of \\
        $\vconsolestream$ and $\literaltostring(\vl)$}.
\end{itemize}

\FormallyParagraph
\begin{mathpar}
\inferrule{
  \env \reverseeqdef (\tenv, \denv, \vconsolestream)\\
  \newenv \eqdef (\tenv, \denv, \vconsolestream \concat \literaltostring(\vl))
}{
  \outputtoconsole(\env, \NVLiteral(\vl)) \evalarrow (\newenv)
}
\end{mathpar}

\SemanticsRuleDef{UnsupportedOutputToConsole}
The function ignores the string value and returns the environment unchanged.

In a runtime without support for a console, the \printstatementsterm{}
in \listingref{semantics-literals} evaluate their list of expressions
with no other effect.

\ProseParagraph
\ProseEqdef{$\newenv$}{$\env$}.

\FormallyParagraph
\begin{mathpar}
\inferrule{}{
  \outputtoconsole(\env, \Ignore) \evalarrow \overname{\env}{\newenv}
}
\end{mathpar}

\section{The Unreachable Statement\label{sec:UnreachableStatement}}
\hypertarget{def-unreachablestatementterm}{}
\listingref{UnreachableStatement} shows an example of using an \unreachablestatementterm{}
to implement a custom form of assertion checking.
\ASLListing{An example use of an \texttt{unreachable} statement}{UnreachableStatement}{\definitiontests/UnreachableStatement.asl}

\subsection{Syntax}
\begin{flalign*}
\Nstmt \derives \ & \Tunreachable \parsesep \Tsemicolon &
\end{flalign*}

\subsection{Abstract Syntax}
\RenderTypes[remove_hypertargets]{stmt_unreachable}


\ASTRuleDef{SUnreachable}

\RenderProseAndFormally{build_stmt_SUnreachable}

\TypingRuleDef{SUnreachable}
\ExampleDef{Typing an Unreachable Statement}
The \unreachablestatementterm{} in \listingref{UnreachableStatement} is well-typed.

\RenderProseAndFormally{annotate_stmt_SUnreachable}

\SemanticsRuleDef{SUnreachable}
\ExampleDef{Evaluating an Unreachable Statement}
Evaluating the specification in \listingref{UnreachableStatement} results in a \dynamicerrorterm,
since the \unreachablestatementterm{} is evaluated.

\RenderProseAndFormally{eval_stmt_SUnreachable}

\section{Pragma Statements\label{sec:PragmaStatements}}
\hypertarget{def-pragmastatementterm}{}

\ASLListing{A pragma statement}{typing-spragma}{\typingtests/TypingRule.SPragma.asl}

\subsection{Syntax}
\begin{flalign*}
\Nstmt \derives \ & \Tpragma \parsesep \Tidentifier \parsesep \ClistZero{\Nexpr} \parsesep \Tsemicolon &
\end{flalign*}

\subsection{Abstract Syntax}
\RenderTypes[remove_hypertargets]{stmt_pragma}


\ASTRuleDef{SPragma}

\RenderProseAndFormally{build_stmt_SPragma}

\TypingRuleDef{SPragma}
\ExampleDef{Typing a Pragma Statement}
The \pragmastatementterm{} in \listingref{typing-spragma} is well-typed.


\RenderProseAndFormally{annotate_stmt_SPragma}
\CodeSubsection{\SPragmaBegin}{\SPragmaEnd}{../Typing.ml}

\subsection{Semantics\label{sec:PragmaSemantics}}
Pragmas are structures present in the \untypedast{} that are designed to be used
by third-party tools.

To avoid conflicts between different ASL parsers, it is recommended that the pragma's identifier $\Tidentifier(\id)$ be prefixed by the name of the ASL tool that supports that pragma
(e.g. ARM for Arm's internal ASL tools). An ASL language processor that does not recognise a pragma directive should generate a warning for that pragma.

Pragmas are not associated with semantics and are discarded from the \typedast.
